% !TeX root = main.tex
% PAGINE ORIGINALI 431-440 — DA COMPILARE
% Incollare qui i quiz del blocco e alla fine le rispettive\registraSoluzione{numero}{lettera}.
% Questo file vuoto non aggiunge pagine al PDF.

% !TeX root = main.tex
% Pagine originali 431-440: 21a simulazione (15 quesiti), 22a simulazione (20 quesiti).

% Pag. originale 431/521
\iffalse Pag. originale 431/521 \fi
\begin{quiz}
$\displaystyle\int_{\pi/2}^{\pi}2\sqrt{1+\cos x}\,dx=$
\begin{enumerate}
\item $-2$
\item $2(\sqrt2-1)$
\item $2$
\item $4(\sqrt2-1)$
\item $4$
\end{enumerate}
\end{quiz}
% SOLUZIONE: D
% SEZIONE_AM1: 5.1
% SPIEGAZIONE: $\sqrt{1+\cos x}=\sqrt2\cos(x/2)$ su $[\pi/2,\pi]$. L\'integrale vale $4\sqrt2[\sin(x/2)]_{\pi/2}^{\pi}=4(\sqrt2-1)$.

% Pag. originale 431/521
\iffalse Pag. originale 431/521 \fi
\begin{quiz}
Il modulo di $z=\dfrac{\sqrt2+i\sqrt2}{\sqrt3-i}$
\begin{enumerate}
\item $\sqrt3$
\item $\sqrt2$
\item $\frac12$
\item $2\sqrt2$
\item $1$
\end{enumerate}
\end{quiz}
% SOLUZIONE: E
% SEZIONE_AM1: 1.5
% SPIEGAZIONE: $|z|=\sqrt{2+2}/\sqrt{3+1}=1$.

% Pag. originale 431/521
\iffalse Pag. originale 431/521 \fi
\begin{quiz}
Lo sviluppo di Mac Laurin di ordine $3$ della funzione $f(x)=\dfrac{e^x-1}{\cos x}$ \`e:
\begin{enumerate}
\item $1+x+\frac{x^2}{2}+\frac23x^3+o(x^3)$
\item $\frac12x^3+o(x^3)$
\item $x+\frac12x^2-\frac13x^3+o(x^3)$
\item $1+x+\frac12x^2-\frac13x^3+o(x^3)$
\item $x+\frac{x^2}{2}+\frac23x^3+o(x^3)$
\end{enumerate}
\end{quiz}
% SOLUZIONE: E
% SEZIONE_AM1: 4.2
% SPIEGAZIONE: $(e^x-1)=x+x^2/2+x^3/6+o(x^3)$ e $1/\cos x=1+x^2/2+o(x^3)$; coefficiente cubico $1/6+1/2=2/3$.

% Pag. originale 431/521
\iffalse Pag. originale 431/521 \fi
\begin{quiz}
La retta tangente alla funzione $f(x)=\dfrac1{1-2e^x}$ in $x=0$ \`e:
\begin{enumerate}
\item $y=-2x$
\item $y=-1-2x$
\item $y=1+2x$
\item $y=2x-1$
\item $y=2x$
\end{enumerate}
\end{quiz}
% SOLUZIONE: D
% SEZIONE_AM1: 4.1
% SPIEGAZIONE: $f(0)=-1$ e $f\prime(0)=2$; la tangente \`e $y=2x-1$.

% Pag. originale 431/521
\iffalse Pag. originale 431/521 \fi
\begin{quiz}
Quanto vale l\'integrale improprio $\displaystyle\frac12\int_0^2\log\left|\frac{x^2}{x-1}\right|\,dx$?
\begin{enumerate}
\item $\log8-2$
\item $2-\log8$
\item $\log4-1$
\item $1-\log4$
\item $+\infty$
\end{enumerate}
\end{quiz}
% SOLUZIONE: C
% SEZIONE_AM1: 5.5
% SPIEGAZIONE: Scrivendo $\log|x^2/(x-1)|=2\log x-\log|x-1|$, l\'integrale completo vale $2\log4-2$, dunque la met\`a \`e $\log4-1$. Le singolarit\`a logaritmiche sono integrabili.

% Pag. originale 432/521
\iffalse Pag. originale 432/521 \fi
\begin{quiz}
L\'asintoto, per $x\to-\infty$, di $f(x)=\dfrac{x^3+x^2}{x^2+1}+\sqrt{x^2+6x}$ \`e:
\begin{enumerate}
\item $y=x+1$
\item $y=-2$
\item $y=-1$
\item $y=2x$
\item nessuna delle altre risposte
\end{enumerate}
\end{quiz}
% SOLUZIONE: B
% SEZIONE_AM1: 3.3
% SPIEGAZIONE: La frazione \`e $x+1+o(1)$; per $x\to-\infty$, $\sqrt{x^2+6x}=-x-3+o(1)$. Quindi $f(x)=-2+o(1)$.

% Pag. originale 432/521
\iffalse Pag. originale 432/521 \fi
\begin{quiz}
Lo sviluppo di Mac Laurin, del terzo ordine, della funzione $f(x)=\sin(\sqrt3\,x)\cdot e^x$ \`e:
\begin{enumerate}
\item $\sqrt3x+\sqrt3x^2+o(x^3)$
\item nessuna delle altre risposte
\item $\sqrt3x+\sqrt3x^2+\frac{3\sqrt3}{2}x^3+o(x^3)$
\item $\sqrt3x-\sqrt3x^2+o(x^3)$
\item $\sqrt3x+\sqrt3x^2+\sqrt3x^3+o(x^3)$
\end{enumerate}
\end{quiz}
% SOLUZIONE: A
% SEZIONE_AM1: 4.2
% SPIEGAZIONE: Il termine cubico \`e $\sqrt3/2-(\sqrt3)^3/6=0$. Pertanto $f(x)=\sqrt3x+\sqrt3x^2+o(x^3)$.

% Pag. originale 432/521
\iffalse Pag. originale 432/521 \fi
\begin{quiz}
Il pi\`u grande intervallo aperto in cui la funzione $f(x)=(x\log x)^8$ \`e decrescente \`e:
\begin{enumerate}
\item $x<e^{-1}\ \lor\ x>1$
\item $(0,e^{-1})\cup(1,+\infty)$
\item $[e^{-1},1]$
\item $(e,+\infty)$
\item $\forall x\in\operatorname{dom}f$
\end{enumerate}
\end{quiz}
% SOLUZIONE: C
% SEZIONE_AM1: 4.3
% SPIEGAZIONE: $f\prime(x)=8(x\log x)^7(\log x+1)$, negativa per $e^{-1}<x<1$. La risposta ufficiale C \`e un intervallo chiuso, bench\'e il testo chieda un intervallo aperto: l\'intervallo aperto corretto sarebbe $(e^{-1},1)$.

% Pag. originale 432/521
\iffalse Pag. originale 432/521 \fi
\begin{quiz}
Sia $z\in\mathbb C$, $z\neq0$. Quale delle seguenti affermazioni NON \`e vera?
\begin{enumerate}
\item $\operatorname{Arg}(z)=-\operatorname{Arg}(\overline z)$
\item $\operatorname{Arg}(z)=\operatorname{Arg}(-z)$
\item $\operatorname{Arg}(z)=\operatorname{Arg}(1/\overline z)$
\item $\operatorname{Arg}(\overline z)=\operatorname{Arg}(1/z)$
\item $\operatorname{Arg}(z\overline z)=0$
\end{enumerate}
\end{quiz}
% SOLUZIONE: B
% SEZIONE_AM1: 1.5
% SPIEGAZIONE: Moltiplicare per $-1$ ruota di $\pi$ l\'argomento: B \`e falsa. Le uguaglianze con $\operatorname{Arg}$ principale richiedono attenzione sui tagli del ramo (per esempio l\'asse reale negativo).

% Pag. originale 432/521
\iffalse Pag. originale 432/521 \fi
\begin{quiz}
$\displaystyle\int_{-1}^{1}\frac{x}{1+x^2}\,dx=$
\begin{enumerate}
\item $-2\log2$
\item $2\arctan1$
\item $0$
\item $2\log2$
\item $\arctan1$
\end{enumerate}
\end{quiz}
% SOLUZIONE: C
% SEZIONE_AM1: 5.1
% SPIEGAZIONE: L\'integranda \`e dispari su un intervallo simmetrico, dunque l\'integrale \`e zero.

% Pag. originale 432/521
\iffalse Pag. originale 432/521 \fi
\begin{quiz}
Sia $f:\mathbb R\to\mathbb R$ una funzione continua e periodica di periodo $2\pi$. Quale delle seguenti affermazioni \`e vera?
\begin{enumerate}
\item $\forall a,b\in\mathbb R:\int_{2a}^{2b}f(x)\,dx=2\int_a^bf(x)\,dx$
\item $\forall a\in\mathbb R:\int_{a-\pi/2}^{a+\pi/2}f(x)\,dx=0$
\item $\int_{-a}^af(x)\,dx=0$
\item $\forall a,b\in\mathbb R:\int_{a+2\pi}^{b+2\pi}f(x)\,dx=\int_a^bf(x)\,dx$
\item nessuna delle altre risposte \`e esatta
\end{enumerate}
\end{quiz}
% SOLUZIONE: D
% SEZIONE_AM1: 5.1
% SPIEGAZIONE: Con la sostituzione $t=x-2\pi$ e la periodicit\`a $f(t+2\pi)=f(t)$, gli integrali in D coincidono. Nell\'originale il limite inferiore di B \`e tipograficamente corrotto.

% Pag. originale 433/521
\iffalse Pag. originale 433/521 \fi
\begin{quiz}
Sia $f$ una funzione due volte derivabile e sia $g(x)=f(\cos x)$. Allora $g\prime\prime(x)=$
\begin{enumerate}
\item $-\cos x\,f\prime(\cos x)-\sin x\,f\prime\prime(\cos x)$
\item $-\cos x\,f\prime(\cos x)$
\item $\sin^2x\,f\prime\prime(\cos x)$
\item $f\prime\prime(-\cos x)\sin^2x-f\prime(-\sin x)\cos x$
\item $f\prime\prime(\cos x)\sin^2x-f\prime(\cos x)\cos x$
\end{enumerate}
\end{quiz}
% SOLUZIONE: E
% SEZIONE_AM1: 4.1
% SPIEGAZIONE: Derivando $g\prime=-\sin x\,f\prime(\cos x)$ con prodotto e catena si ottiene E.

% Pag. originale 433/521
\iffalse Pag. originale 433/521 \fi
\begin{quiz}
L\'insieme dei numeri reali $\alpha$ per cui $\displaystyle\lim_{x\to0^+}\frac{x-\sin x}{x^\alpha}>0$ \`e
\begin{enumerate}
\item $\alpha\leq3$
\item $\alpha\geq2$
\item $\alpha>2$
\item $\alpha>3$
\item $\alpha\geq3$
\end{enumerate}
\end{quiz}
% SOLUZIONE: E
% SEZIONE_AM1: 3.4
% SPIEGAZIONE: $x-\sin x\sim x^3/6$, dunque il limite \`e $0$ per $\alpha<3$, $1/6$ per $\alpha=3$ e $+\infty$ per $\alpha>3$.

% Pag. originale 433/521
\iffalse Pag. originale 433/521 \fi
\begin{quiz}
Quale delle seguenti propriet\`a NON \`e una propriet\`a dell\'integrale definito su un intervallo limitato $I$?
\begin{enumerate}
\item $\int_a^bf(x)\,dx=\int_a^cf(x)\,dx+\int_c^bf(x)\,dx$, $a,b,c\in I$
\item $\int_a^b(\alpha f(x)+\beta g(x))\,dx=\alpha\int_a^bf(x)\,dx+\beta\int_a^bg(x)\,dx$, $a,b\in I$, $\alpha,\beta\in\mathbb R$
\item $\int_a^bf(x)\,dx\geq0$, $a,b\in I$
\item se $f\leq g$ $\forall x\in I$, $a<b$, allora $\int_a^bf(x)\,dx\leq\int_a^bg(x)\,dx$
\item $\int_a^bf(x)\,dx\leq\int_a^bf(x)\,dx\leq\int_a^b|f(x)|\,dx$, $a,b\in I$, $a<b$
\end{enumerate}
\end{quiz}
% SOLUZIONE: C
% SEZIONE_AM1: 5.1
% SPIEGAZIONE: L\'integrale di una funzione generica pu\`o essere negativo; C non vale senza $f\geq0$ e $a<b$.

% Pag. originale 433/521
\iffalse Pag. originale 433/521 \fi
\begin{quiz}
L\'area della regione finita di piano delimitata dalla $f(x)=x(1-x)$ e dalla retta di equazione $g(x)=-\frac x2$ vale
\begin{enumerate}
\item $\frac5{12}$
\item $\frac1{12}$
\item $-\frac1{12}$
\item $\frac9{16}$
\item nessuna delle altre risposte \`e esatta
\end{enumerate}
\end{quiz}
% SOLUZIONE: D
% SEZIONE_AM1: 5.1
% SPIEGAZIONE: Le intersezioni sono $x=0$ e $x=3/2$. L\'area \`e $\int_0^{3/2}(3x/2-x^2)\,dx=9/16$.

% Pag. originale 436/521
\iffalse Pag. originale 436/521 \fi
\begin{quiz}
La funzione $f:[a,e]\to\mathbb R$, $f(x)=\log|x|$, $a\leq e$, soddisfa le ipotesi del teorema degli zeri se:
\begin{enumerate}
\item $a>0$
\item $a\in(0,+\infty)$
\item $a\in(0,e)$
\item $a\in(0,1)$
\item $a\in(0,1]$
\end{enumerate}
\end{quiz}
% SOLUZIONE: D
% SEZIONE_AM1: 4.3
% SPIEGAZIONE: Occorrono continuit\`a su $[a,e]$ (quindi $a>0$) e segni opposti agli estremi: $\log a<0<\log e=1$, cio\`e $0<a<1$.

% Pag. originale 436/521
\iffalse Pag. originale 436/521 \fi
\begin{quiz}
La parte principale, per $x\to+\infty$, della funzione $f(x)=2x-\sqrt[3]{x^3+x^4}$ \`e
\begin{enumerate}
\item $x-x^{4/3}$
\item $2x+o(x)$
\item $2x$
\item $-x^{4/3}$
\item $-\frac13x^3$
\end{enumerate}
\end{quiz}
% SOLUZIONE: D
% SEZIONE_AM1: 3.4
% SPIEGAZIONE: $\sqrt[3]{x^3+x^4}\sim x^{4/3}$ domina $2x$, quindi $f(x)\sim-x^{4/3}$.

% Pag. originale 436/521
\iffalse Pag. originale 436/521 \fi
\begin{quiz}
La successione $\{a_n=3n+\sin n,\ n\in\mathbb N\}$
\begin{enumerate}
\item \`e limitata
\item $\lim_{n\to+\infty}a_n=-\infty$
\item $\lim_{n\to+\infty}a_n=+\infty$
\item $\lim_{n\to+\infty}a_n$ non esiste
\item $\lim_{n\to+\infty}a_n=\pm\infty$
\end{enumerate}
\end{quiz}
% SOLUZIONE: C
% SEZIONE_AM1: 2.2
% SPIEGAZIONE: $3n-1\leq a_n\leq3n+1$, quindi $a_n\to+\infty$.

% Pag. originale 436/521
\iffalse Pag. originale 436/521 \fi
\begin{quiz}
$\displaystyle\lim_{x\to1}\frac{\log x^2}{\log(1+\log x)}=$
\begin{enumerate}
\item $0$
\item $\log2$
\item $1$
\item $2$
\item $+\infty$
\end{enumerate}
\end{quiz}
% SOLUZIONE: D
% SEZIONE_AM1: 3.1
% SPIEGAZIONE: $\log x^2=2\log x$ vicino a $1$ e $\log(1+\log x)\sim\log x$; rapporto $2$.

% Pag. originale 436/521
\iffalse Pag. originale 436/521 \fi
\begin{quiz}
Sia data la funzione $f(x)=\begin{cases}1-x&x\leq1\\ e^x-e&x>1.\end{cases}$
\begin{enumerate}
\item $f([1,8])$ \`e un sottoinsieme proprio di $f([0,8])$
\item $f(( -\infty,0])=f([0,+\infty))$
\item $f(x)<0$ per qualche $x\in\mathbb R$
\item $f([1,8])=f([0,8])$
\item $f([0,8])$ \`e un sottoinsieme proprio di $f([1,8])$
\end{enumerate}
\end{quiz}
% SOLUZIONE: D
% SEZIONE_AM1: 4.3
% SPIEGAZIONE: Su $[0,1]$ l\'immagine \`e $[0,1]$ e su $[1,8]$ \`e $[0,e^8-e]$; quest\'ultimo contiene $[0,1]$, quindi le immagini coincidono.

% Pag. originale 436/521
\iffalse Pag. originale 436/521 \fi
\begin{quiz}
Se $f(0)=1$, $f(5)=10$ e se $f$ \`e continua, allora sicuramente esiste una $x_0\in(0,5)$ tale che $f(x_0)=$
\begin{enumerate}
\item $2$
\item $0$
\item $5e$
\item $\frac12$
\item $e-2$
\end{enumerate}
\end{quiz}
% SOLUZIONE: A
% SEZIONE_AM1: 4.3
% SPIEGAZIONE: Il teorema dei valori intermedi garantisce ogni valore strettamente compreso fra $1$ e $10$, in particolare $2$.

% Pag. originale 437/521
\iffalse Pag. originale 437/521 \fi
\begin{quiz}
La derivata di $f(x)=\log(e^{4x}\cos x)$ \`e
\begin{enumerate}
\item $\frac{\sin x}{\cos x}$
\item $4-\tan x$
\item $4+\tan x$
\item $4\tan x$
\item $e^{4x}\sin x$
\end{enumerate}
\end{quiz}
% SOLUZIONE: B
% SEZIONE_AM1: 4.1
% SPIEGAZIONE: Nel dominio, $f(x)=4x+\log(\cos x)$; quindi $f\prime(x)=4-\tan x$.

% Pag. originale 437/521
\iffalse Pag. originale 437/521 \fi
\begin{quiz}
Qual \`e una primitiva della funzione $f(x)=e^x\sin(3e^x)$?
\begin{enumerate}
\item $\frac13\cos(3e^x)$
\item $-\frac13\cos(3e^x)$
\item $e^x\sin(3e^x)+3e^x\cos(3e^x)$
\item $\cos(3e^x)$
\item $\frac13\cos(e^x)$
\end{enumerate}
\end{quiz}
% SOLUZIONE: B
% SEZIONE_AM1: 5.1
% SPIEGAZIONE: Ponendo $t=3e^x$ si ha $dt=3e^x\,dx$ e $\int e^x\sin(3e^x)\,dx=-\frac13\cos(3e^x)+C$.

% Pag. originale 437/521
\iffalse Pag. originale 437/521 \fi
\begin{quiz}
L\'integrale improprio $\displaystyle\int_0^{+\infty}\frac{x\sqrt{x}+\sin x}{x^2+1}\,dx$
\begin{enumerate}
\item converge
\item diverge positivamente
\item \`e indeterminato
\item diverge negativamente
\item non \`e possibile stabilire il comportamento
\end{enumerate}
\end{quiz}
% SOLUZIONE: B
% SEZIONE_AM1: 5.5
% SPIEGAZIONE: Per $x\to+\infty$, il numeratore \`e asintotico a $x^{3/2}$ e il denominatore a $x^2$: integranda $\sim1/\sqrt{x}$, il cui integrale diverge positivamente.

% Pag. originale 437/521
\iffalse Pag. originale 437/521 \fi
\begin{quiz}
L\'equazione differenziale $y\prime=\dfrac{x^2+y}{x^2+1}$
\begin{enumerate}
\item \`e lineare
\item \`e a variabili separabili
\item ha soluzioni costanti
\item l\'unica soluzione \`e $y(0)=1$
\item nessuna delle risposte precedenti
\end{enumerate}
\end{quiz}
% SOLUZIONE: A
% SEZIONE_AM1: 6.1
% SPIEGAZIONE: $y\prime-\frac{1}{x^2+1}y=\frac{x^2}{x^2+1}$ \`e lineare del primo ordine.

% Pag. originale 437/521
\iffalse Pag. originale 437/521 \fi
\begin{quiz}
Sapendo che lo sviluppo di Mac Laurin \`e $f(x)=-1+2x+3x^2+o(x^2)$, allora
\begin{enumerate}
\item $f\prime(x)=2+6x+o(x)$ nell\'intorno di $x=0$
\item $x=0$ \`e punto di massimo
\item $x=0$ \`e punto di minimo
\item $f\prime(x)=2+3x+o(x)$ nell\'intorno di $x=0$
\item $f\prime(x)=1+6x+o(x)$ nell\'intorno di $x=0$
\end{enumerate}
\end{quiz}
% SOLUZIONE: A
% SEZIONE_AM1: 4.2
% SPIEGAZIONE: La risposta ufficiale A deriva dalla derivazione formale dello sviluppo. In generale derivare un resto $o(x^2)$ non garantisce un resto $o(x)$ senza ipotesi aggiuntive: quesito matematicamente non rigoroso.

% Pag. originale 437/521
\iffalse Pag. originale 437/521 \fi
\begin{quiz}
$\displaystyle\lim_{x\to0}\frac{\sinh^2(5x)}{1-\cos(3x)}=$
\begin{enumerate}
\item $\frac{25}{3}$
\item $\frac53$
\item $\frac{50}{3}$
\item $+\infty$
\item $\frac{50}{9}$
\end{enumerate}
\end{quiz}
% SOLUZIONE: E
% SEZIONE_AM1: 3.1
% SPIEGAZIONE: $\sinh^2(5x)\sim25x^2$, $1-\cos(3x)\sim9x^2/2$, rapporto $50/9$.

% Pag. originale 438/521
\iffalse Pag. originale 438/521 \fi
\begin{quiz}
Sia $z=(k-3)+i(2-k)$; $w=z^2$ \`e reale
\begin{enumerate}
\item se $k=-3$ oppure $k=-2$
\item se $k=0$ oppure $k=1$
\item per nessun valore di $k$
\item $\forall k\in\mathbb R$
\item se $k=3$ oppure $k=2$
\end{enumerate}
\end{quiz}
% SOLUZIONE: E
% SEZIONE_AM1: 1.5
% SPIEGAZIONE: La parte immaginaria di $z^2$ \`e $2(k-3)(2-k)$, nulla per $k=2$ o $k=3$.

% Pag. originale 438/521
\iffalse Pag. originale 438/521 \fi
\begin{quiz}
Lo sviluppo di Mac Laurin, del $3^\circ$ ordine, della funzione $f(x)=\sinh\left(x-\frac{x^3}{6}\right)$ \`e
\begin{enumerate}
\item $x$
\item $x-\frac{x^3}{6}$
\item $x+o(x)$
\item $x+o(x^3)$
\item $x+o(x^2)$
\end{enumerate}
\end{quiz}
% SOLUZIONE: D
% SEZIONE_AM1: 4.2
% SPIEGAZIONE: $\sinh t=t+t^3/6+o(t^3)$ con $t=x-x^3/6$; i termini cubici si cancellano. Anche C ed E sono affermazioni vere ma meno precise: ambiguit\`a della fonte.

% Pag. originale 438/521
\iffalse Pag. originale 438/521 \fi
\begin{quiz}
L\'equazione differenziale $y\prime=\cos(y-x)+1$
\begin{enumerate}
\item \`e lineare del I ordine
\item \`e a variabili separabili
\item ammette $y(x)=x$ come soluzione particolare
\item $\cos(y-x)=-1$ \`e la soluzione costante
\item nessuna delle altre soluzioni \`e esatta
\end{enumerate}
\end{quiz}
% SOLUZIONE: E
% SEZIONE_AM1: 6.1
% SPIEGAZIONE: Non \`e lineare n\'e separabile in $x,y$; $y=x$ darebbe $1=2$, falso. L\'opzione D \`e una condizione algebrica, non una soluzione costante.

% Pag. originale 438/521
\iffalse Pag. originale 438/521 \fi
\begin{quiz}
Sia $f$ una funzione definita su tutto $\mathbb R$, a valori positivi, per la quale esiste $\lim_{x\to+\infty}f(x)$. Allora
\begin{enumerate}
\item $\lim_{x\to+\infty}f(x)\geq0$ e pu\`o essere $+\infty$
\item $\lim_{x\to+\infty}f(x)>0$ e pu\`o essere $+\infty$
\item $\lim_{x\to+\infty}f(x)=+\infty$
\item $\lim_{x\to+\infty}f(x)=\ell$, $\ell\in[0,+\infty)$
\item $\lim_{x\to+\infty}f(x)=\ell$, $\ell\in\mathbb R$
\end{enumerate}
\end{quiz}
% SOLUZIONE: A
% SEZIONE_AM1: 3.1
% SPIEGAZIONE: Un limite di funzione positiva pu\`o essere zero, positivo finito oppure $+\infty$.

% Pag. originale 438/521
\iffalse Pag. originale 438/521 \fi
\begin{quiz}
Sia $f$ una funzione integrabile su ogni intervallo $[0,b]\subseteq\mathbb R$ e tale che $f(x)\geq0$, $\forall x\geq0$. Allora
\begin{enumerate}
\item l\'integrale improprio $\int_0^{+\infty}f(x)\,dx$ pu\`o essere indeterminato
\item l\'integrale improprio $\int_0^{+\infty}f(x)\,dx$ necessariamente converge
\item l\'integrale improprio $\int_0^{+\infty}f(x)\,dx$ necessariamente diverge positivamente
\item $\lim_{c\to+\infty}\int_0^cf(x)\,dx=\ell$, $\ell\in(0,+\infty)$ oppure $+\infty$
\item $\lim_{c\to+\infty}\int_0^cf(x)\,dx=\ell$, $\ell\in[0,+\infty)$
\end{enumerate}
\end{quiz}
% SOLUZIONE: D
% SEZIONE_AM1: 5.5
% SPIEGAZIONE: Le primitive integrali parziali sono non decrescenti e hanno limite in $[0,+\infty]$. La risposta ufficiale D esclude impropriamente il caso $f\equiv0$, per cui il limite \`e $0$: nessuna alternativa descrive esattamente tutti i casi.

% Pag. originale 438/521
\iffalse Pag. originale 438/521 \fi
\begin{quiz}
Supponiamo che $\forall\varepsilon>0\ \exists\delta>0$ tale che $\forall |x|<\delta\Rightarrow |f(x)|<\varepsilon x^2$, allora necessariamente
\begin{enumerate}
\item $f(x)=o(x^2)$, $x\to0$
\item $f(x)\sim x^2$, $x\to0$
\item $f(x)\sim x^3$, $x\to0$
\item $x^2=o(f(x))$, $x\to0$
\item nessuna delle altre risposte \`e esatta
\end{enumerate}
\end{quiz}
% SOLUZIONE: A
% SEZIONE_AM1: 3.4
% SPIEGAZIONE: \`E esattamente la definizione di $f(x)=o(x^2)$ per $x\to0$ (per $x\ne0$).

% Pag. originale 439/521
\iffalse Pag. originale 439/521 \fi
\begin{quiz}
$\displaystyle\lim_{x\to+\infty}\frac3x\log(e^{2x}+x)=$
\begin{enumerate}
\item $1$
\item $3$
\item $+\infty$
\item $6$
\item $2$
\end{enumerate}
\end{quiz}
% SOLUZIONE: D
% SEZIONE_AM1: 3.1
% SPIEGAZIONE: $\log(e^{2x}+x)=2x+\log(1+xe^{-2x})=2x+o(1)$; il prodotto tende a $6$.

% Pag. originale 439/521
\iffalse Pag. originale 439/521 \fi
\begin{quiz}
Siano date le successioni $\{a_n\},\{b_n\}$ con $a_n=b_n+(-1)^n\dfrac{n}{n^2+1}$. Allora necessariamente
\begin{enumerate}
\item la convergenza di $\{a_n\}$ non implica la convergenza di $\{b_n\}$
\item $\{a_n\}$ pu\`o convergere indipendentemente dal comportamento di $\{b_n\}$
\item se $\{a_n\}$ converge anche $\{b_n\}$ deve convergere
\item se $\{b_n\}$ \`e indeterminata allora non si pu\`o dire quale sia il comportamento di $\{a_n\}$
\item nessuna delle altre risposte \`e esatta
\end{enumerate}
\end{quiz}
% SOLUZIONE: C
% SEZIONE_AM1: 2.2
% SPIEGAZIONE: $a_n-b_n=(-1)^nn/(n^2+1)\to0$, quindi le due successioni hanno lo stesso limite quando esiste.
Risposta completata





